\section{Introduction and results}\label{intro:main}
Non-isolated flat connections are a basic obstruction to interpreting a semiclassical expansion as a sum of ordinary Gaussian saddle contributions. In a plumbing assembled from simple torus pieces, this issue already occurs in a sharp form: closed invariants can have leading logarithms, and the logarithmic order can grow with the number of pieces. The main question here is how classical critical geometry, absolute weight, and phase fit together in an actual finite quantum sum.

We restrict to connected finite untwisted trees with at least two vertices, with \(+2,-2\) caps at every vertex and the orientation-reversing section--fibre exchange on each backbone edge. The primitive section is the base boundary curve of the punctured-sphere product block. The twist/framing conventions agree with the earlier capped-tree calculations \cite{P48,P49,P49B,P4849}. We use \(r=k+2\geq 3\), \(S_{ab}=\sqrt{2/r}\sin(\pi ab/r)\), \(Z_r(S^2\times S^1)=1\), and \(\tau_r=Z_r/S_{11}\), so \(\tau_r(S^3)=1\). For a fibre colour \(a\), its trace is \(2\cos(\pi a/r)\), and its negative-centre deficit is \((r-a)/r\). Geometric cutoffs and Wilson insertions always refer to this specified marking and ribbon framing.

\paragraph{The classical family and its scale faces.}
The exact flat relations force noncentral fibres to form an independent set. They also give a combinatorial description of the maximal-dimensional actual solution branches. The vertices exposed by some maximum matching and their neighbours are removed first; the remaining forest is perfectly matched. This residual forest controls the number of leading logarithms. Nonzero scaled adjacent charge products belong to an ambient near-flat lift, rather than to the exact flat space itself. The geometric description retains both these scales and the conjugacy information that radial traces alone forget.

\paragraph{The exact invariant by action.}
For the finite modular contraction, the signed angular alias sums reduce to finite rational-cone sums. The forest incidence matrices provide the integrality needed for uniform expansions across intersections of additive walls. Cap-core insertions then yield bounded phase filters and identify the two actual Chern--Simons actions, \(0\) and \(1/2\) modulo \(1\). Their expansion coefficients are distributions in the original fibre traces. Each phase has a power--logarithm expansion without any further level-residue oscillation.

\paragraph{The absolute comparison.}
The reduced classical action in the relevant irreducible charts is obtained from the flat critical equations and the cyclic comparison of the finite deformation complexes; the latter is needed to identify the function, not merely its critical set. Primitive cohomology bases, the stabilizer, and the transverse Gaussian fix the normalization independently of the finite sum. The leading degeneracy localizes the matched angle dots at zero. On the quantum side the actual low-rank ribbon transfer, including its signed tails, then yields the same conditional joint law for the supported Wilson routes, including words connecting distinct dormant components. At fixed transition products the signed kernel is essential. Independent positive radial cutoffs would change the coefficient.

The paper keeps the two levels of conclusion separate: absolute classical weights and joint dormant-angle comparisons at leading order, together with selected transition laws; complete higher radial expansion by action with supported coefficient distributions. The latter leaves a further problem of assigning higher coefficients to individual angular components with the same radial data. This distinction is used in the statements and examples throughout.


\paragraph{Guide to the proof.}
There are three calculations with different roles. Sections~\ref{geo:sec} and~\ref{classic:sec:action} concern actual flat connections and the reduced classical action. The radial and phase sections start from finite modular sums and determine the signed expansions by action. The angular and Hall sections compare a stated class of genuine ribbon probes to the independently normalized classical model. The leading comparison does not use a conjectured expansion of a continuum functional integral.
